\begin{corollary}\label{cor:divisorial-current-ordinary-type}
Let $S=\sum_jc_j[D_j]$ be a closed positive divisorial $(1,1)$-current of finite mass on a connected compact complex manifold. It determines a canonical ordinary quasi-plurisubharmonic type with representatives $v$ satisfying
\[
 \ddc v=S-\beta
\]
for a smooth closed real form $\beta$. Any two such representatives differ smoothly. For every $a\in\QPSH$, this type is characterized by
\[
 a\preceq v\quad\Longleftrightarrow\quad
 a\preceq_{\mathcal I}v\quad\Longleftrightarrow\quad
 \nu(a,D_j)\geq c_j\ \hbox{for all }j.
\]
For every Hermitian background $\Omega$ admitting the chosen representative,
\[
 P_\Omega[v]=Q_\Omega[v]\simeq v.
\]
In particular the Gauduchon bound in \cref{thm:trace-divisorial-current-budget} gives this canonical ordinary type for $S_Z(u;Y)$ on each fixed compact model $Z$.
\end{corollary}
\begin{proof}
\Cref{lem:bc-one-one-current-description} writes $S=\beta+\ddc v$ with $\beta$ smooth and closed and $v$ quasi-plurisubharmonic. If another pair is used, the difference of the potentials has smooth $\ddc$. Taking a smooth Hermitian trace and applying local elliptic regularity shows that the difference is smooth. Thus the ordinary type depends only on $S$.

The three-way characterization is \cref{lem:divisorial-potential-domination}. The ideal envelope preserves all multiplier ideals, so its generic Lelong number on each $D_j$ is still $c_j$. The same lemma gives $Q_\Omega[v]\preceq v$. Also $Q_\Omega[v]\geq v-\sup v$, so these functions are ordinarily equivalent. Consequently $Q_\Omega[v]$ is admissible in the ordinary envelope of $v$, giving $Q_\Omega[v]\leq P_\Omega[v]$. The reverse inequality is \cref{prop:p-below-i}. This proves the envelope identity and ordinary equivalence. The final assertion uses the positive closed current of finite mass already constructed from the Gauduchon estimate.
\end{proof}

This realization describes the divisorial part on one specified model. It does not identify the higher-codimensional singularities needed to recover every $\tau_{Y,u}(E)$ on further modifications. In the curve-fibration setting, \cref{thm:curve-fibration-exact-trace} gives that additional identification. For general trace data it remains a separate problem. The canonical ordinary type here is selected by the smooth residual term; other potentials with the same multiplier ideals need not belong to it.
